\documentclass[a4paper, 12pt, oneside]{article}
\usepackage{array}
\usepackage{etaremune}
\usepackage{amsmath}
\usepackage{multirow}
\usepackage{setspace}
\usepackage{booktabs}
\usepackage{environ}
\usepackage{graphicx}
\usepackage{biblatex}
\usepackage{geometry}
\usepackage{caption}
\usepackage{wrapfig}
\usepackage[table]{xcolor}
\usepackage{fancyhdr}


\begin{document}
	\thispagestyle{empty}
	{\raggedleft
	{\large \textbf{ITTA Special HAZMAT Publication 800-001A}}\\\rule{\textwidth}{3pt} \-
	{\huge \textbf{Guideline for using \\
	HAZMAT respirator regulation 800-001A}} \- \\ \rule[0pt]{\textwidth}{3pt}
	\begin{figure}[b]
		\caption*{\rule{\textwidth}{2pt}}
		\includegraphics[width=2cm, height=2cm]{C:/Users/akos/Downloads/SRA_logo}
		\tiny \caption*{| v2.0.0 |}
	\end{figure}
	
	\newpage
	
	\centering \large \textbf{Authority} \\\raggedright \normalsize
	This publication for the Ruby Mountains Superstructure has been developed by the ITAA in coordinance with the SRA. This guideline is responsible for ensuring proper production and safety standards, including minimum requirements for such equipment, however, this standard may be overridden given proper clearance to do so.
	
	Nothing in this publication shall be released to the general public, as it is confidential. \\ \-\vfill 
	\centering \large Institution for Technological Action and Advocacy \\
	Special HAZMAT Publication 800-001A \\ \- \\\normalsize
	\begin{tabular}{>{\columncolor{lightgray}}c}
		Certain commercial entities, equipment, or materials may be identified in this document \\ 
		in order to describe an experimental procedure or concept adequately. \\ 
		Such identification is not intended to imply recommendation or endorsement by ITAA, \\
		nor is it intended to imply that the entities, materials, or equipment are \\
		necessarily the best available for the purpose.\\
	\end{tabular} \\ \-\\
	\textbf{Comments on this publication may be submitted to:} \\
	the Board of Directors \\ 
	The Administrator (with permission)\\ \-\\
	
	\newpage
	\pagestyle{fancy}
	
	\fancyhead[L]{ITAA SP HAZ 800-001A} \-\\
	\centering \large \textbf{Abstract} \\\raggedright \normalsize \-\\
	This regulation document is intended to serve hazmat units within the Superstructure, and to adapt to newer technologies now available, along with increased relative safety standards. This Special Publication (SP) is the first one of its kind, but will be iterated on.\\ \-\\\mbox{ }
	
	\centering \large \textbf{Keywords} \\\raggedright \normalsize \-\\
	confidentiality; critical infrastructure; Administrative directives; mandates; policy; standards. \\ \-\\\mbox{ }
	
	\centering \large \textbf{Acknowledgments} \\\raggedright \normalsize \-\\
	The author wishes to thank those colleagues who ensured this regulation didn't sink miles below standard. The author also wishes to  gratefully acknowledge and appreciate the many comments from the public and private sectors whose 'thoughtful' and 'constructive' comments improved the quality and usefulness of this publication.
	
	\newpage
	
	\tableofcontents
	
	\newpage
	
	\section{INTRODUCTION}	
	\subsection{Background and Purpose}
	HAZMAT publications of the Institute for Technological Action and Advocacy (ITAA) provide guidance regarding equipment usage, manufacturing requirements, situation control, and what to use in said situations. \\
	This document provides basic guidance on how to identify what filter should be used in HAZMAT situations where SCBA equipment is either unavailable or not required, along with other information regarding filter construction, proper usage, and traits. \\
	\subsection{Acronyms}
	\begin{tabular}{|p{2cm} p{10cm}}
		ITAA & Institution for Technological Action and Advocacy \\
		SRA & Superstructure Research Authority \\
		BoD & Board of Directors \\
		HAZMAT & Hazardous-Materials \\
		SCBA & Self-Contained Breathing Apparatus \\
		ppm & parts per million \\
		SP & Special Publication \\
		v/v & volume per volume \\
		IS & Inserts Chart \\
		IEC & Insert Efficiency Chart \\
		SD & Short Duration \\
		MD & Medium Duration \\
		HD & High Duration \\
		XHD & Extra-High Duration \\
		Hi & High \\
		Lo & Low \\
	\end{tabular}
		\subsection{Document Organization}
			- Section 1 provides an introduction to this document, including its background and purpose, and a list of acronyms used herein. \\
			- Section 2 provides information on basic filter identification. \\
			- Section 3 describes information on what should be the base main filter types, the size standards of said filter types, and accepted filter casing materials. \\
			- Section 4 provides other information.
			
	\newpage
			
	\section{Filter Identification}
		This section provides basic filter identification, along with a list \\
		of the main filter insert types.
	\subsection{Filter Trait Charts}
	\subsubsection{Inserts Chart}
		Each insert can be applied into a filter, be it one or more.
		\begin{tabular}{|p{1cm}||p{10cm}|p{1.57cm}|}
			\hline
			Types & Traits & Colors \\\hline
			AX & Gases and vapours of organic compounds with a boiling point below 65 degrees. & \cellcolor[HTML]{995000} \\
			A & Gases and vapours of organic compounds with a boiling point above 65 degrees. & \cellcolor[HTML]{997000} \\
			B & Inorganic gases \& vapours$^{[1]}$. & \cellcolor[HTML]{CACACA} \\
			E & Sulphur dioxide \& hydrogen chloride. & \cellcolor[HTML]{CCBB00} \\
			K & Ammonia and its organic derivatives. & \cellcolor[HTML]{00AA00} \\
			CO & Carbon monoxide. & \cellcolor[HTML]{000000} \\
			Hg & Mercury vapour. & \cellcolor[HTML]{990000} \\
			NO & Nitrous gases, including nitrogen monoxide. & \cellcolor[HTML]{0000BB} \\
			RK & Radioactive iodine, including radioactive methyl iodine. & \cellcolor[HTML]{E96000} \\
			P & Particulates.$^{[2]}$ & \cellcolor[HTML]{424242} \\\hline
		\end{tabular} \-\\
		\subsubsection{Insert Efficiency Chart}
			Filtration efficiency ranges from 1-5, 5 being the best and 1 the worst. \\
			Each filter insert can have its own filtering efficiency, \\
			be it to reduce complexity / cost or conserve space.
			\begin{tabular}{|p{1cm}||p{4cm}|p{7.57cm}|}
				\hline
				Types & Reduction (ppm) & Reduction (Particulate, ppm) \\\hline
				C5 & 90,000 ppm (9\% v/v) & 999,999 ppm (99.9999\% v/v) - (log-6) \\
				C4 & 50,000 ppm (5\% v/v) & 900,000 ppm (90\% v/v) - (log-1) \\
				C3 & 10,000 ppm (1\% v/v) & 800,000 ppm (80\% v/v) - (~log-0.7) \\
				C2 & 5,000  ppm (0.5\% v/v) & 700,000 ppm (70\% v/v) - (~log-0.53) \\
				C1 & 1,000  ppm (0.1\% v/v) & 600,000 ppm (60\% v/v) - (~log-0.4) 	\\\hline
			\end{tabular} \\ \- \\
			$[1]$: Chlorine, Hydrogen sulphide, etc. \\
			$[2]$: Particulate filter reduction amount ranges from
			 5 - ($^{99.9999\%}_{log-6red.}$) to 1 - ($^{60\%}_{0.4}$)
			
	\newpage
			
	\section{Main filter types} \-\\
		\subsection{Main Filter Types}
			\subsubsection{Main Filter Duration Chart*}
				\begin{tabular}{|p{1cm}||p{8.44cm}|p{3cm}|}
					\hline
					Types & Durations & Avg. Durations \\\hline
					SD & 1H(High-exposure) - 8H(Low-exposure) &  ~4.5 Hours\\
					MD & 4H(High-exposure) - 18H(Low-exposure) & ~11 Hours \\
					HD & 8H(High-exposure) - 30H(Low-exposure) & ~19 Hours \\
					XHD & 12H(High-exposure) - 60H(Low-exposure) & ~36 Hours \\\hline
				\end{tabular} \\ \-\\
			\subsubsection{Main Filter Type Chart}
				\begin{tabular}{|p{3cm}||p{8.44cm}|p{1cm}|}
					\hline
					Types & Traits & Types \\\hline
					Omni(Combined) & Must contain all inserts. & Hi/Lo \\
					REACTOR & Must contain RK-type insert. & Hi/Lo \\
					Particulate & P-type and B-type inserts only. & Hi/Lo \\
					Vapour & May contain all inserts excluding RK-type. & Hi/Lo \\\hline	
				\end{tabular} \\ \-\\
				\small *Do note that these are the expected durations, and can be exceeded, however if the design falls short of these standards, the BoD may be advised in order to test if the filter is still sufficient for use. \normalsize \\ \-\\
		\subsection{Main Filter Parameters}
			Large filters (above diameterx10) are to be connected to a hose connecting to both the respirator and filter, to be stored inside a suitable containter. \\ \-\\
			
			\subsubsection{Sizes \& Casings Chart}
				\begin{tabular}{|p{4cm}||p{3cm}|p{5.45cm}|}
					\hline
					Types & Sizes & Accpeted Casing Materials \\\hline
					Hi-Omni[SD] & 12(cm)x10(cm) & Nomex / Titanium Blend \\
					Lo-Omni[SD] & 10(cm)x10(cm) & Nomex / Titanium Blend \\
					Hi-Omni[MD] & 12(cm)x15(cm) & Nomex / Titanium Blend \\
					Lo-Omni[MD] & 10(cm)x15(cm) & Nomex / Titanium Blend \\
					Hi-Omni[HD] & 12(cm)x20(cm) & Nomex / Titanium Blend \\
					Lo-Omni[HD] & 10(cm)x20(cm) & Nomex / Titanium Blend \\
					Hi-Omni[XHD] & 12(cm)x30(cm) & Nomex / Titanium Blend \\
					Lo-Omni[XHD] & 10(cm)x30(cm) & Nomex / Titanium Blend \\
					Hi-REACTOR[MD] & 12(cm)x10(cm) & Nomex / Titanium Blend \\
					Lo-REACTOR[MD] & 10(cm)x10(cm) & Nomex / Titanium Blend \\
					Hi-REACTOR[HD] & 12(cm)x15(cm) & Nomex / Titanium Blend \\
					Lo-REACTOR[HD] & 10(cm)x15(cm) & Nomex / Titanium Blend \\
					Hi-REACTOR[XHD] & 12(cm)x25(cm) & Nomex / Titanium Blend \\
					Lo-REACTOR[XHD] & 10(cm)x25(cm) & Nomex / Titanium Blend \\
					Hi-Particulate[MD] & 12(cm)x5(cm) & Nomex \\
					Lo-Particulate[MD] & 10(cm)x5(cm) & Nomex \\
					Hi-Particulate[HD] & 12(cm)x10(cm) & Nomex \\
					Lo-Particulate[HD] & 10(cm)x10(cm) & Nomex \\
					Hi-Particulate[XHD] & 12(cm)x15(cm) & Nomex \\
					Lo-Particulate[XHD] & 10(cm)x15(cm) & Nomex \\
					Hi-Vapour[SD] & 12(cm)x5(cm) & Nomex / Titanium Blend \\
					Lo-Vapour[SD] & 10(cm)x5(cm) & Nomex / Titanium Blend \\
					Hi-Vapour[MD] & 12(cm)x10(cm) & Nomex / Titanium Blend \\
					Lo-Vapour[MD] & 10(cm)x10(cm) & Nomex / Titanium Blend \\
					Hi-Vapour[HD] & 12(cm)x15(cm) & Nomex / Titanium Blend \\
					Lo-Vapour[HD] & 10(cm)x15(cm) & Nomex / Titanium Blend \\
					Hi-Vapour[XHD] & 12(cm)x25(cm) & Nomex / Titanium Blend \\
					Lo-Vapour[XHD] & 10(cm)x25(cm) & Nomex / Titanium Blend \\\hline
				\end{tabular} \-\\
				\small Format is as following: Radius(cm/mm) x length(cm). \\
				Other filter casing sizes \& materials can be discussed with the BoD.
			
	\newpage
			
	\section{Other}\-\\
		\subsection{Copyright}
		These fictional fan documents of the White Knuckle universe may not be:
		redistributed, changed or used for commercially without explicit permission 
		from the author.
		
		© 2026 moron\_notification, All rights reserved.
		\subsection{Sources}
			\subsubsection{Links}
				\url{https://nvlpubs.nist.gov/nistpubs/SpecialPublications/NIST.SP.800-175A.pdf}i (helped in document construction, also borrowed some text to reformat)
				\url{https://whiteknuckle.miraheze.org/wiki/White_Knuckle_Wiki} (bery goodies game, go play it, this is a fan-document for wuckle anyways) \\ \-\\
			\subsubsection{People}
				lizuwubeth (xmatyo5) | [Critic] \\
				Wold (wold123) | [Moral Help] \\
				JB-20 (joyousbicycle20) | [Quality Assurance] \\
				gkugfk3 (gkugfk3) | [Quality Assurance] \\
	\begin{figure}[b]
		\centering
		\includegraphics[width=5cm, height=5cm]{C:/Users/akos/Downloads/itaa-gud}
		\caption*{TAKE YOUR COGNITIVOL}
	\end{figure}
	
\end{document}